\chapter{可验证性设计}\label{ch:verif}

本章说明设计为验证提供的观察手段。验证平台、用例、核对项、测量口径和结果见测试报告。

\section{设计原则}

设计让每条消息和每一帧的去向都能只从顶层端口确定，从而满足三点：
\begin{itemize}
  \item \textbf{只看端口。}系统级验证不需要探测内部信号，同一套激励可以直接用于 RTL 和综合后网表。
  \item \textbf{期望值可独立生成。}帧格式、CRC、校验和、BAG 间隔和 IC/RM 判决都按协议规则定义，参考模型可以不读 RTL 独立算出期望结果。
  \item \textbf{逐条核对。}每条消息、每一帧恰好对应一个结果脉冲，验证可以逐条比对去向，而不只是比较总数。
\end{itemize}

\section{可观察性}

\begin{itemize}
  \item \textbf{事件端口。}\code{ev_tx}、\code{ev_a}、\code{ev_b}、\code{ev_rm_drop}、\code{ev_deliver} 对每条发送消息给出 accept 或四种丢弃之一，对线上每一帧在其所在网络给出一个类别脉冲，完好帧再给出一个结果脉冲（第~\ref{sec:rxclass}~节）。脉冲出现的时间有确定上界：发送消息的脉冲在末字节后 96 拍内，接收类别脉冲在帧尾后 32 拍内，结果脉冲在帧尾后 150 拍内。
  \item \textbf{计数器。}27 个统计计数器经配置口读回，读值的取样时刻有精确定义（$t$ 拍发出读请求，$t+2$ 拍返回 $t-1$ 拍及以前的脉冲数，第~\ref{ch:cfg}~章），可以与事件记录逐拍对齐。
  \item \textbf{线路字节。}GMII 输出直接来自 MAC 的输出寄存器，帧的每个字节和起始拍都在端口上可见；两网同拍起始，便于对照两网的输出。
  \item \textbf{服务状态。}RTC 的偏差和校时次数、SNMP 的请求和 community 错误计数、615A 的状态、原因码和加载区内容都可经配置口读回（附录~\ref{app:regs}）。
  \item \textbf{RAM 碰撞监视。}\code{afdx_ram_sdp} 在仿真中检查同拍读写同一地址，出现即报错；该监视用 \code{`ifndef SYNTHESIS} 包住，不影响综合网表。它直接检验第~\ref{sec:cfgok}~节“避让”设计是否成立。
\end{itemize}

\section{可控性}

链路上的各类异常都可以从端口构造：在 GMII 输入上改动字节、丢弃或重放整帧，即可产生 FCS 错、帧头错、序号乱序、双网时间差等情况；配置口可以随时改写 VL 配置以构造配置不允许、未知 VL 和队列满。设计不需要专门的测试模式。
